\Cref{tab:2-position} summarises the preceding methods. Neither code
size nor search organisation alone establishes superiority: graph
traversal, compressed estimation and two-level refinement spend memory
and candidate work differently. At nine bits per dimension, IVF-RaBitQ
matches JHQ's one-bit primary plus eight-bit residual payload; four- and
eight-bit variants are also tested. Chapters 4--6 investigate JHQ's
construction, selection and budget choices. Chapter 7 compares achieved
recall, total memory and query cost under Chapter 3's controls.

\begin{table}[!htbp]
\centering
\caption[Related approaches and JHQ-GPU design questions]{Related approaches and the design questions they motivate for JHQ-GPU.}
\label{tab:2-position}
\footnotesize\setlength{\tabcolsep}{4pt}
\begin{tabularx}{\linewidth}{@{}>{\raggedright\arraybackslash}p{.20\linewidth}*{2}{>{\raggedright\arraybackslash}X}@{}}
\toprule
Approach & Main trade-off & Question for JHQ-GPU \\
\midrule
IVF-PQ / OPQ & Codebooks are trained for low reconstruction error; coordinates within a subspace may be coupled & How can the primary code's coordinate structure simplify GPU scoring? \\
Fast Scan / Quicker ADC & Small tables accelerate lookup under hardware and precision constraints & Does splitting tables repay extra GPU lookups? \\
SONG / CAGRA & Graph navigation reduces scoring but adds graph storage and dependent access & Where can a compressed scan compete with resident graph search? \\
BANG & Host storage extends capacity at the cost of communication & What can a complete compressed index resident on one GPU achieve? \\
IVF-RaBitQ & Multi-bit codes trade storage and probe effort against distance-estimation quality & How does JHQ's two-stage search compare with IVF-RaBitQ at matched recall? \\
Adaptive effort & Quality control may require labels, prediction or calibration & Can workload agreement select a useful budget at acceptable quality and cost? \\
\bottomrule
\end{tabularx}
\end{table}
